\documentclass[10pt,a4paper]{amsart}
\usepackage[margin=19mm]{geometry}
\usepackage{amsmath,amssymb,mathtools}
\usepackage[scr=rsfs,cal=euler]{mathalfa}
\usepackage{xfrac}
\usepackage[unicode,hidelinks,bookmarks=false]{hyperref}
\newcommand{\ie}{i.e.\ }
\newcommand{\cf}{cf.\ }
\newcommand{\resp}{respectively\ }
\newcommand{\Subsection}{\subsection}
\providecommand{\nobreakdash}{\nobreak\hbox{-}}
\setlength{\emergencystretch}{2em}
\multlinegap0pt
\usepackage{tikz}
\usepackage{amssymb}
\usepackage{stackengine}
\usepackage{enumerate}
\usepackage{xfrac}
\newtheorem{theorem}{Theorem}
\title{Isoclinic groups and conjugacy quandles}
\author{Mohamad Maassarani}
\date{Source: arXiv:2606.22554v1 (CC BY 4.0)}
\begin{document}
\maketitle

\noindent\emph{Source and licence.} Isoclinic groups and conjugacy quandles. Version arXiv:2606.22554v1 (CC BY 4.0), distributed by arXiv under the Creative Commons Attribution 4.0 International licence, http://creativecommons.org/licenses/by/4.0/, as recorded in the arXiv licence metadata for this exact version and shown on the abstract page https://arxiv.org/abs/2606.22554v1. Full licence notice: \texttt{licenses/arxiv-2606.22554-CC-BY-4.0.txt}.\par\smallskip

\noindent\textit{Editorial scope.} Counted unit: the article's theorem theorem (source0.tex lines 124--126). The article's complete original proof of it is delivered with it (source0.tex lines 127--166, one \texttt{proof} environment of 40 lines). No mathematics is written, repaired or re-derived: the statement and the proof are the article's own lines, in the article's own notation, and no retained line is substituted or re-flowed.  No further source-local context is needed: every cross-reference of every retained line resolves inside the two delivered spans.  Nothing was omitted from any retained span, so the package carries no omission record.  No retained line contains a citation, so the wrapper carries no bibliography and no bibliography entry.  No retained line contains a figure environment, an image-inclusion command or drawing code, so no image is delivered and the package declares no asset dependency.  Editorial formatting only, in addition: an AMS \texttt{amsart} preamble that restates the article's own declarations and macros used by the retained lines, namely the environments \texttt{theorem} on separate counters, together with the standard packages those lines need.  The retained environments print in this document as Theorem 1 in order of appearance, whereas the article prints the same items with the numbering of its own full document.  The complete native source of the article as distributed in the arXiv e-print for this exact version is bundled unchanged as \texttt{sources/203357/source0.tex}.
\begin{theorem}
If $G$ and $H$ are two isoclinc finite groups of the same order then their conjugacy quandles $Conj(G)$ and $Conj(H)$ are isomorphic.
\end{theorem}

\begin{proof}
Let $G$ and $H$ be as in the theorem. We hence have two compatible isomorphisms $\alpha :G/Z(G)\to H/Z(H)$ and $\phi : G' \to H'$ 
andthe map $\alpha$ induces an isomorphism on the abelianisation $\bar{\alpha} : G/Z(G)^{ab} \to H/Z(H)^{ab}$. Let $s_G$ be a section to the morphism $G\to G/Z(G)^{ab}$ and $s_H$ be a section to $H\to H/Z(H)^{ab}$ as in $3)$ of the previous lemma. By $1)$ of the previous lemma and the previous proposition the distinct elements of $G$ and $H$ are the elements :
$g_igs_G(x)$ and $h_ihs_H(y)$ for $i\in\{1,\dots,n\}$ $g\in G'$, $h\in H'$, $ x\in G/Z(G)^{ab}$ and $y\in H/Z(H)^{ab}$ and where $n=n_G=n_H$ and $g_1,\dots,g_n $ are some elements of $Z(G)$ and $h_1,\dots,h_n$ are some elements in $Z(H)$. We hence have a bijection $\psi : G\to H$ biven by :
$$\psi(g_igs_G(x))=h_i\phi(g)s_H(\bar{\alpha}(x)),$$
for $i\in \{1,\dots,n\}$, $g\in G'$ and $x\in G/Z(G)^{ab}$. We will prove that $\psi$ is a quandle morphism, i.e $\psi(aba^{-1})=\psi(a)\psi(b) \psi(a)^{-1}$ for $a,b\in G$. Take two elements $g_igs_G(x)$ and $g_jg's_G(x')$ of $G$ as above. We have that : 
\begin{equation*}
\begin{split}
\psi(g_igs_G(x) g_jg's_G(x')(g_igs_G(x))^{-1})&=\psi([g_igs_G(x),g_jg's_G(x')]_Gg_jg's_G(x'))\\
&= \psi(g_j([g_igs_G(x),g_jg's_G(x')]_Gg')s_G(x'))\\
&=h_j\phi([g_igs_G(x),g_jg's_G(x')]_Gg')s_H(\bar{\alpha}(x'))\\
&=h_j\phi([gs_G(x),g's_G(x')]_Gg')s_H(\bar{\alpha}(x'))\\
&=h_j\phi([gs_G(x),g's_G(x')]_G)\phi(g')s_H(\bar{\alpha}(x'))
\end{split}
\end{equation*}
On the other hand :
\begin{equation*}
\begin{split}
\psi(g_igs_G(x))\psi(g_jg's_G(x'))\psi(g_igs_G(x))^{-1}&=h_i\phi(g)s_H(\bar{\alpha}(x))h_j\phi(g')s_H(\bar{\alpha}(x'))(h_i\phi(g)s_H(\bar{\alpha}(x)))^{-1} \\
&=[h_i\phi(g)s_H(\bar{\alpha}(x)),h_j\phi(g')s_H(\bar{\alpha}(x'))]_Hh_j\phi(g')s_H(\bar{\alpha}(x'))\\
&=h_j[\phi(g)s_H(\bar{\alpha}(x)),\phi(g')s_H(\bar{\alpha}(x'))]_H\phi(g')s_H(\bar{\alpha}(x'))
\end{split}
\end{equation*}
Compairing the results of the two last equations we get that $\psi$ is a quandle morphism if and only if :
\begin{equation}\label{eq1}
\phi([gs_G(x),g's_G(x')]_G)=[\phi(g)s_H(\bar{\alpha}(x)),\phi(g')s_H(\bar{\alpha}(x'))]_H
\end{equation}
Using the compatibility of $\phi$ and $\alpha$, the fact that we have chosen $s_H$ such that $\alpha(\overline{s_G(x)})=\overline{s_H(\bar{\alpha}(x))}$ for all $x\in G/Z(G)^{ab}$, then applying $4)$ of the last lemma we get :
\begin{equation}
\begin{split}
\phi([gs_G(x),g's_G(x')]_G)&=[\alpha(\overline{gs_G(x)}),\alpha(\overline{g's_G(x')})]_H \\
&= [\alpha(\overline{g}) \alpha(\overline{s_G(x)}),\alpha(\overline{g'}) \alpha(\overline{s_G(x')})]_H \\
&=[\alpha(\overline{g}) \overline{s_H(\bar{\alpha}(x)}),\alpha(\overline{g'})\overline{s_H(\bar{\alpha}(x'))}]_H \\
&=[\overline{\phi(g)} \: \overline{s_H(\bar{\alpha}(x)}),\overline{\phi(g')}\:\overline{s_H(\bar{\alpha}(x'))}]_H\\
&=[\overline{\phi(g)s_H(\bar{\alpha}(x)}),\overline{\phi(g')s_H(\bar{\alpha}(x'))}]_H\\
&=[\phi(g)s_H(\bar{\alpha}(x)),\phi(g')s_H(\bar{\alpha}(x'))]_H
\end{split}
\end{equation}
We have proved that equation (\ref{eq1}) holds and hence $\psi$ is a quandle morphism. This proves the theorem since we have seen that $\psi$ is a bijection.
\end{proof} 

\end{document}
